% エントリ 23(2026-11-15): 2 世代目の BH が生まれたときのスピンと質量の分布
\section{[E23, 2026-11-15] 2 世代目の BH が生まれたときのスピンと質量の分布}\label{sec:e23}

2 世代目の BH は合体の残骸で、スピン 0.69 のまわりで生まれる。この研究では、それを自然スピンの仮定に頼らない試験体に使う(DESIGN.md D2)。これまで、GW241011 の主星の生まれたスピンは一つの値(0.72 か 0.78)に固定していた(E9--E22)。E10 で残骸のスピンの式から分布を出したが、親の集団にも、ボソンにも、解析にもつないでいなかった。ここで、解析が呼べる部品として書き、親の集団と年齢の模型(E22)につないだ(ユーザー指示「スピン分布を書く」、2026-11-14)。

結論を先に書く。
\begin{enumerate}
\item \textbf{部品}: \texttt{share/remnant/remnant.py}(numpy だけ)。親の対の質量と、合体のときのスピンから、残骸の質量・スピン・キックを返す。スピンは \arxiv{1605.01938} の式 16、質量は \arxiv{1206.3803} の式 18、キックは \arxiv{1605.01067} の式 44--47。$10^6$ 対で 0.2--0.3 秒。
\item \textbf{分布(ボソンなし)}: 基準の組(親のスピンは Beta(0.69, 2.28)、向きは等方、質量比は 0.5--1 の一様)で、中央値 0.675、5--95\% は 0.597--0.758。0.69 以上で生まれるのは 33\%、0.72 以上は 15\%、0.78 以上は 2.5\% である。質量は親の質量の和の 0.954 倍。
\item \textbf{0.78 は裾にある}: GW241011 の主星の測定(0.78、90\% 区間 0.69--0.87)の中心より上で生まれる 2 世代目は 40 個に 1 個で、区間に入るのも 3 個に 1 個である。これまでの「0.78 で生まれた」は楽観的な端にあたる。
\item \textbf{ボソンは親にも効く}: 親($7$--$14\Msun$)が合体までに超放射で回転を失うと、2 世代目の生まれたスピンの上の裾が削られる。$f\ge10^{14}\GeV$ では 0.72 以上が 15\% から 3\% に減り、0.78 以上はなくなる。窓の $f=10^{12}\GeV$ では、0.78 以上が 2.5\% から 1.5\%(1 世代目の年齢が古い組なら 0.6\%)になる。
\item \textbf{生まれたスピンと年齢の両方について周辺化した GW241011}: 0.69 以上にとどまる確率がボソンなしの値の半分になる $f$ は、2 世代目の年齢が「広い」組で $4.2\times10^{11}$、「星団の目安」で $9.1\times10^{11}$、「遅い」で $1.4\times10^{11}\GeV$。生まれたスピンを 0.78 に固定した値(E22)の 0.49--0.57 倍、0.72 に固定した値の 0.81--0.93 倍で、0.71--0.72 で生まれた場合にあたる。E22 の「同じ $f$ で両立する度合い」は、いちばん有利な組で、$20\Msun$ の 1 世代目が 0.88 から 0.72、$10\Msun$ が 0.33 から 0.12 に下がる。
\item \textbf{星団に残る条件}: キックが脱出速度より小さい残骸だけを使うと、分布は狭くなり(脱出速度 100 km/s で残るのは 13\%、0.72 以上は 7\%)、$f$ はさらに 0.76--0.84 倍になる。ただしキックの式は粗く、球状星団の模擬計算の値の一つを再現できなかった。既定では選ばない。
\end{enumerate}

計算は \texttt{calc/e23\_natal2g.py}(出力 \texttt{calc/out/e23.json})。$\chi_\mathrm{SR}$ は単体版(E21 の版)、年齢は E22 の模型、$\mu=0.71\times10^{-12}\eV$。

\subsection{定義と模型}

\paragraph{残骸の式}
親を 1(重いほう)、2 と書き、$q=m_2/m_1\le1$、$\nu=q/(1+q)^2$、スピンの大きさを $a_1$、$a_2$ とする。$\beta$、$\gamma$ はスピン 1、2 と軌道角運動量 $\boldsymbol L$ のなす角、$\alpha$ はスピンどうしのなす角である(この節の $\alpha$ は角で、重力の結合定数ではない)。

残骸のスピン $\chi_b$ は、Hofmann, Barausse \& Rezzolla(\arxiv{1605.01938})の式 13--16 で
\begin{align}
a_\mathrm{tot}&=\frac{a_1\cos\beta+a_2\cos\gamma\,q^2}{(1+q)^2},\qquad a_\mathrm{eff}=a_\mathrm{tot}+\xi\nu\,(a_1\cos\beta+a_2\cos\gamma),\nonumber\\
|\boldsymbol\ell|&=\Bigl|L_\mathrm{ISCO}(a_\mathrm{eff})-2a_\mathrm{tot}\bigl[E_\mathrm{ISCO}(a_\mathrm{eff})-1\bigr]+\sum_{i=0}^{1}\sum_{j=0}^{2}k_{ij}\,\nu^{1+i}a_\mathrm{eff}^{\,j}\Bigr|,\nonumber\\
\chi_b&=\frac{1}{(1+q)^2}\Bigl[a_1^2+a_2^2q^4+2a_1a_2q^2\cos\alpha+2\bigl(a_1\cos\beta+a_2q^2\cos\gamma\bigr)|\boldsymbol\ell|\,q+|\boldsymbol\ell|^2q^2\Bigr]^{1/2}.
\label{eq:e23spin}
\end{align}
$E_\mathrm{ISCO}$、$L_\mathrm{ISCO}$ は Kerr 時空の最内安定円軌道のエネルギーと角運動量(同 式 2--6)、係数は同 Table 1 の上段($\xi=0.41616$)で、$k_{00}$ は等質量・無スピンで $\chi_b=0.68646$ になる条件(同 式 11)から決める。E10 の \texttt{hbr16.py} と同じ式である。

残骸の質量 $M_b$ は、Barausse, Morozova \& Rezzolla(\arxiv{1206.3803})の式 18 で
\begin{equation}
\frac{M_b}{m_1+m_2}=1-\bigl[1-E_\mathrm{ISCO}(a_\mathrm{tot})\bigr]\nu-4\nu^2\bigl[4p_0+16p_1a_\mathrm{tot}(a_\mathrm{tot}+1)+E_\mathrm{ISCO}(a_\mathrm{tot})-1\bigr],
\label{eq:e23mass}
\end{equation}
$p_0=0.04827$、$p_1=0.01707$(同 式 12)。等質量・無スピンで 0.952 倍になる。

キック(反跳)の速さ $v_k$ は、Gerosa \& Kesden(\arxiv{1605.01067})の式 44--47 を使う。スピンの組み合わせ(同 式 36)
\begin{equation}
\boldsymbol\Delta=\frac{a_1\hat{\boldsymbol S}_1-q\,a_2\hat{\boldsymbol S}_2}{1+q},\qquad\tilde{\boldsymbol\chi}=\frac{a_1\hat{\boldsymbol S}_1+q^2a_2\hat{\boldsymbol S}_2}{(1+q)^2}
\label{eq:e23delta}
\end{equation}
($\hat{\boldsymbol S}_{1,2}$ はスピンの向きの単位ベクトル)の、$\boldsymbol L$ に平行な成分($\parallel$)と垂直な成分の大きさ($\perp$)で
\begin{align}
v_k&=\sqrt{v_m^2+2v_mv_\perp\cos\zeta+v_\perp^2+v_\parallel^2},\qquad v_m=A\nu^2\frac{1-q}{1+q}(1+B\nu),\qquad v_\perp=H\nu^2\Delta_\parallel,\nonumber\\
v_\parallel&=16\nu^2\Bigl[\Delta_\perp\bigl(V_{11}+2V_A\tilde\chi_\parallel+4V_B\tilde\chi_\parallel^2+8V_C\tilde\chi_\parallel^3\bigr)+2\tilde\chi_\perp\Delta_\parallel\bigl(C_2+2C_3\tilde\chi_\parallel\bigr)\Bigr]\cos\Theta
\label{eq:e23kick}
\end{align}
と書く。$A=1.2\times10^4$、$H=6.9\times10^3$、$V_{11}=3677.76$、$V_A=2481.21$、$V_B=1792.45$、$V_C=1506.52$、$C_2=1140$、$C_3=2481$(どれも km/s)、$B=-0.93$、$\zeta=145^\circ$。$\Theta$ は合体のときの落下の向きに関わる角で、同論文にならって $[0,\pi]$ の一様にとる。

\paragraph{$\boldsymbol\Delta$ の符号}
式~\eqref{eq:e23delta} の $\boldsymbol\Delta$ には、重いほうのスピンが正の符号で入る。この符号は結果に効く。$\cos\zeta<0$ なので、重いほうが $\boldsymbol L$ と同じ向きに回っている($\Delta_\parallel>0$)と、スピンの項が質量比の項を一部打ち消す。$q=0.5$ で重いほうだけが 0.5 のスピンを持つとき、軌道面内のキックは、順行で 91、無スピンで 157、逆行で 258 km/s である。同論文の公開コード(v1.0.3)は同じ符号で、注記に「v1.0.2 で負号を足した。論文の誤植」とある。arXiv の本文(ar5iv で読んだ版)の式 36 もこの符号で、脚注に、ほかの文献には符号の食い違いがあると書かれている。ここでは直したあとの符号を使い、逆の符号にした場合も下で比べる。

\paragraph{親の集団と、ボソンの入り方}
親のスピン $a_{1,2}$ は\textbf{合体のときの値}である。ボソンがなければ、生まれたときの値(1 世代目の自然スピン)に等しい。ボソンがあれば、親は生まれてから合体するまで(1 世代目の年齢 $T_1$)に超放射で回転を失うので、$\chi_\mathrm{SR}$ を通した値になる。2 世代目は、こうして決まった $M_b$ と $\chi_b$ で生まれ、次の合体まで(2 世代目の年齢 $T_2$)にもう一度 $\chi_\mathrm{SR}$ を通る。
\begin{equation}
(m_i,\,\chi_{i,0})\ \xrightarrow{\ \chi_\mathrm{SR},\ T_1\ }\ (m_i,\,a_i)\ \xrightarrow{\ \text{式 \eqref{eq:e23spin}、\eqref{eq:e23mass}}\ }\ (M_b,\,\chi_b)\ \xrightarrow{\ \chi_\mathrm{SR},\ T_2\ }\ \chi_\mathrm{2G}.
\label{eq:e23chain}
\end{equation}
親の集団についての仮定は表~\ref{tab:e23par} のとおり。

\begin{table}[H]
\centering\footnotesize
\begin{tabular}{lll>{\raggedright\arraybackslash}p{6.0cm}}
\toprule
量 & 基準の組 & 動かした範囲 & 根拠\\
\midrule
親の質量の和 & $20.6\Msun$ & --- & 残骸が GW241011 の主星の $19.6\Msun$ になる\\
親の質量比 & 0.5--1 の一様 & 1、0.25--1 の一様 & 星団の中の合体は等質量に寄る(\arxiv{1906.10260} Sec.~IV B)\\
親の生まれたスピン & Beta(0.69, 2.28) & Beta(2, 5)、0、0.2、0.5 & 1 世代目の推定(\arxiv{2607.01317} の事後分布の中央値。E12)\\
スピンの向き & 等方 & --- & 星団の中で対になった BH\\
親の年齢 & 1 世代目の既定(E22) & 古い組、赤方偏移 0.05、1 & 親の合体の赤方偏移は 0.3 にとる\\
星団に残る条件 & なし & $v_k<50$、100、300 km/s & キックの式が粗い(下)\\
\bottomrule
\end{tabular}
\caption{親(1 世代目)の集団の仮定。}
\label{tab:e23par}
\end{table}

\subsection{部品と確かめ}

\texttt{remnant.py} の \texttt{merge} は、親の対の配列を受けとり、スピンの向きを等方に引いて、残骸の質量、スピン、キックを返す。確かめたことは次のとおり(\texttt{e23\_natal2g.py check})。
\begin{itemize}
\item スピンの式は、E10 の \texttt{hbr16.py} と 40 万個の入力で同じ値を返す(差は 0)。等質量・無スピンで 0.68646、$q=1/3$・無スピンで 0.540、等質量で 0.5 ずつ揃ったスピンで 0.832、質量比 $\to0$ で重いほうのスピンに戻る。$k_{00}=-3.821$(原文は $\approx-3.82$)。
\item 質量の式は、等質量・無スピンで 0.95173 倍、等質量で最大のスピンが揃ったときの放出が 9.95\%(原文 Sec.~2 の値)、質量比 $\to0$ で $[1-E_\mathrm{ISCO}(0)]\nu$ になる。
\item キックは、無スピンでは $q=0.362$ で最大の 175.2 km/s(\arxiv{1906.10260} Sec.~IV D の「$q\sim1/3$ で $175\pm11$ km/s」)、等質量・無スピンで 0、等質量で最大のスピンが軌道面内で反平行のとき $V_{11}$、最大のスピンで向きを乱数で探した最大は 4912 km/s(原文は最大を約 5000 km/s としている)。
\item \texttt{merge} は、親をどちらの順で渡しても同じ値を返す。引いた向きの余弦の平均は $10^{-3}$ 以下。
\item スピンの当てはめが 1 を越えるのは、重いほうのスピンが 0.99 を越えて $\boldsymbol L$ に揃い、質量比が 0.03 付近のときだけで、最大 0.0014。\texttt{merge} は 1 で切る。
\item 渡した先で照らす自己点検を付けた(\texttt{python remnant.py}。決まった入力での 14 個の値と、値の範囲)。
\end{itemize}

\subsection{生まれたときのスピンと質量(ボソンなし)}

40 万対から作った分布を表~\ref{tab:e23dist} と図~\ref{fig:e23}(a) に示す。

\begin{table}[H]
\centering\small
\begin{tabular}{lllcccccc}
\toprule
親のスピン & 質量比 & $v_\mathrm{esc}$ [km/s] & 残る割合 & 中央値 & 5--95\% & $\ge0.69$ & $\ge0.72$ & $\ge0.78$\\
\midrule
Beta(0.69, 2.28) & 0.5--1 & なし & 1.000 & 0.675 & 0.597--0.758 & 0.331 & 0.146 & 0.025\\
 &  & 50 & 0.034 & 0.686 & 0.657--0.718 & 0.302 & 0.046 & 0.003\\
 &  & 100 & 0.130 & 0.683 & 0.644--0.728 & 0.294 & 0.071 & 0.006\\
 &  & 300 & 0.575 & 0.672 & 0.605--0.733 & 0.249 & 0.080 & 0.009\\
 & 1 & なし & 1.000 & 0.689 & 0.627--0.762 & 0.486 & 0.199 & 0.024\\
 & 0.25--1 & なし & 1.000 & 0.657 & 0.495--0.753 & 0.249 & 0.116 & 0.023\\
\midrule
Beta(2, 5) & 0.5--1 & なし & 1.000 & 0.677 & 0.591--0.757 & 0.388 & 0.178 & 0.018\\
 &  & 100 & 0.075 & 0.686 & 0.626--0.752 & 0.454 & 0.179 & 0.013\\
0 & 0.5--1 & なし & 1.000 & 0.675 & 0.631--0.686 & 0 & 0 & 0\\
 &  & 100 & 0.623 & 0.682 & 0.669--0.686 & 0 & 0 & 0\\
0.2 & 0.5--1 & なし & 1.000 & 0.675 & 0.609--0.721 & 0.318 & 0.057 & 0\\
 &  & 100 & 0.115 & 0.690 & 0.638--0.730 & 0.499 & 0.116 & 0\\
0.5 & 0.5--1 & なし & 1.000 & 0.693 & 0.554--0.793 & 0.517 & 0.366 & 0.089\\
 &  & 100 & 0.033 & 0.716 & 0.572--0.808 & 0.617 & 0.484 & 0.175\\
\bottomrule
\end{tabular}
\caption{2 世代目の生まれたときのスピン $\chi_b$ の分布(ボソンなし、親の質量の和 $20.6\Msun$、向きは等方)。$v_\mathrm{esc}$ の行は、キックが $v_\mathrm{esc}$ より小さい残骸だけの分布。右の 3 列は、その値以上で生まれる割合。}
\label{tab:e23dist}
\end{table}

\begin{itemize}
\item \textbf{基準の組}: 中央値 0.675、5--95\% は 0.597--0.758。山は等質量・無スピンの値 0.686 のすぐ下にある。親の質量比が 1 より小さいと下がり、親のスピンが両側に広げる。最大は 0.913 で、0.87 以上は 0.04\% である。
\item \textbf{親の質量比}: 等質量だけにすると中央値 0.689、0.25--1 の一様にすると 0.657 で、0.69 以上の割合は 0.49 から 0.25 まで動く。0.78 以上の割合は 2.3--2.5\% でほとんど変わらない。
\item \textbf{親のスピン}: 親が回らなければ 0.631--0.686 に収まり、0.69 以上は一つもない。上の裾は親のスピンの分布の裾で決まる。Beta(2, 5) では 0.78 以上が 1.8\%、親がみな 0.5 なら 8.9\% である。E10 の値(Beta(2, 5) で中央値 0.677、95 パーセンタイル 0.757)を再現する。
\item \textbf{角度の補正}(\arxiv{1605.01938} の式 18。当てはめの偏り $-0.005$ を $-0.001$ にする)を入れると、基準の組の中央値は 0.001 下がり、0.69 以上は 0.333 から 0.324、0.72 以上は 0.148 から 0.143 になる。
\item \textbf{質量}: 基準の組で、親の質量の和の 0.954 倍(5--95\% は 0.945--0.962 倍)。親の和を目標の質量の $1/0.952$ 倍にとれば、残骸は目標の $\pm1\%$ に入る。
\end{itemize}

\paragraph{GW241011 の主星との関係}
測定は $\chi_1=0.78$、90\% 区間は 0.69--0.87 である(E10)。基準の組で、生まれたときにこの区間に入る 2 世代目は 33\%、中心の 0.78 より上は 2.5\% である。E10 の「典型的な 2 世代目は、区間の下端のすぐ外にある」を数にするとこうなる。生まれたスピンを 0.78 に固定した計算(E8--E22 の楽観的な側)は、40 個に 1 個の場合を代表にしていたことになる。

\subsection{星団に残る条件と、文献の数値実験との照合}

2 世代目がもう一度合体するには、キックを受けたあとも星団に残っていなければならない。式~\eqref{eq:e23kick} のキックは、基準の組で中央値 249 km/s(5--95\% は 61--1011 km/s)で、脱出速度 $v_\mathrm{esc}=50$、100、300 km/s より小さいものは 3.4\%、13.0\%、57.5\% である。
\begin{itemize}
\item \textbf{基準の組では、残るのは親のスピンが小さい対である}。分布は狭くなる。$v_\mathrm{esc}=100$ km/s で 5--95\% は 0.644--0.728、0.72 以上は 15\% から 7\%、0.78 以上は 2.5\% から 0.6\% に減る。0.69 以上の割合は 33\% から 29\% で、あまり変わらない。
\item \textbf{親のスピンが一つの値なら、残るものは上にずれる}。親がみな 0.2 のとき、0.69 以上は 32\% から 50\% に増え、みな 0.5 のとき中央値は 0.693 から 0.716 に上がる。重いほうが順行のときにキックが小さいためである。数値相対論の代理模型による研究(\arxiv{2503.21278})の「残った残骸の山は 0.69 より上に来る」と同じ向きである。
\end{itemize}

\paragraph{球状星団の模擬計算との照合}
\arxiv{1906.10260}(CMC)は、1 世代目のスピンを一つの値に置いた星団の模擬計算で、残骸が星団に残る割合(Sec.~III A)と、もう一度合体した 2 世代目のスピン(Sec.~IV D)を示している。同じ親のスピンで比べた(質量比は 0.5--1 の一様、200 万対)。

\begin{table}[H]
\centering\small
\begin{tabular}{lccc}
\toprule
 & 残る割合 & 中央値 & 5--95\%\\
\midrule
親が回らない、$v_\mathrm{esc}=50$ km/s & 0.324 & 0.685 & 0.683--0.686\\
親が回らない、$v_\mathrm{esc}=100$ km/s & 0.623 & 0.682 & 0.669--0.686\\
\arxiv{1906.10260}、親が回らない & 約 0.6 & --- & 0.67--0.69\\
\midrule
親のスピン 0.5、$v_\mathrm{esc}=50$ km/s & 0.006 & 0.711 & 0.565--0.815\\
親のスピン 0.5、$v_\mathrm{esc}=100$ km/s & 0.034 & 0.718 & 0.572--0.808\\
同、$\boldsymbol\Delta$ の符号を逆、$v_\mathrm{esc}=50$ km/s & 0.006 & 0.679 & 0.550--0.808\\
同、$\boldsymbol\Delta$ の符号を逆、$v_\mathrm{esc}=100$ km/s & 0.039 & 0.660 & 0.533--0.800\\
\arxiv{1906.10260}、親のスピン 0.5 & 0.03 未満 & 0.62 & 0.53--0.72\\
\bottomrule
\end{tabular}
\caption{残った残骸のスピン。文献の行は、同論文の本文の値(スピンは、もう一度合体した 2 世代目の 90\% の範囲)。}
\label{tab:e23cmc}
\end{table}

\begin{itemize}
\item \textbf{親が回らない場合は合う}。$v_\mathrm{esc}=100$ km/s で残る割合は 62\%、残ったもののスピンは 0.669--0.686 である。
\item \textbf{親のスピンが 0.5 の場合は、残る割合は合うが、スピンが合わない}。残ったものの中央値は 0.71--0.72 で、文献の 0.62 より 0.1 高い。$\boldsymbol\Delta$ の符号を逆にすると(公開コードが直す前の符号)、中央値は 0.66--0.68 に下がり、下の端も 0.53--0.55 に下がるが、上の端は 0.80--0.81 のままで文献の 0.72 に届かない。質量比を 1 や 0.25--1 に変えても、中央値は 0.70--0.72 である。
\item \textbf{原因は決められなかった}。文献の値は星団の力学を通ったあとの「もう一度合体した 2 世代目」のもので、質量比と脱出速度の分布がこの比較と違う。残骸の当てはめの式も別のもので、合体のときのスピンの向きの決め方は本文に書かれていない。
\end{itemize}

この食い違いは、キックで選んだ分布をどこまで信じるかに関わる。残る割合が数 \% 以下の標本は、キックがたまたま小さい特別な配置で、式~\eqref{eq:e23kick} の近似(二つの項に同じ位相 $\cos\Theta$ を掛けていること、$\boldsymbol\Delta$ の符号)にいちばん敏感な所である。そのうえ \arxiv{2503.21278}(Sec.~2.3)は、この系統の式のキックが、数値相対論の代理模型より平均で 10 倍大きいとしている。そこで、\textbf{既定ではキックで選ばず}、選んだ場合($v_\mathrm{esc}=50$--300 km/s)は頑健さの確認に使うことにする。基準の組では、残るのは親のスピンが小さい対なので(残る割合 13\%)、この弱点の効き方は、親がみな 0.5 の場合より小さい。

\subsection{親が超放射で回転を失う場合}

親の質量は $6.9$--$13.7\Msun$ で、$\mu=0.71\times10^{-12}\eV$ では重力の結合定数が 0.036--0.073、211 の臨界スピン($4\alpha_g/(1+4\alpha_g^2)$。$\alpha_g$ は結合定数)が 0.15--0.29 である。それより上で生まれた親は、時間があれば臨界まで回転を失う。親の年齢を E22 の 1 世代目の模型から引き、$\chi_\mathrm{SR}$ を通して式~\eqref{eq:e23chain} の前半を計算した(20 万対、表~\ref{tab:e23boson})。

\begin{table}[H]
\centering\small
\begin{tabular}{lcccc@{\qquad}cccc}
\toprule
 & \multicolumn{4}{c}{1 世代目の年齢: 既定} & \multicolumn{4}{c}{1 世代目の年齢: 古い}\\
$f$ [GeV] & 親の平均 & $\ge0.69$ & $\ge0.72$ & $\ge0.78$ & 親の平均 & $\ge0.69$ & $\ge0.72$ & $\ge0.78$\\
\midrule
ボソンなし & 0.232 & 0.331 & 0.146 & 0.025 & 0.232 & 0.331 & 0.146 & 0.025\\
$10^{11.5}$ & 0.230 & 0.330 & 0.144 & 0.023 & 0.228 & 0.328 & 0.142 & 0.021\\
$10^{12}$ & 0.220 & 0.322 & 0.131 & 0.015 & 0.207 & 0.313 & 0.114 & 0.006\\
$10^{12.5}$ & 0.192 & 0.301 & 0.093 & 0.007 & 0.162 & 0.279 & 0.052 & 0.0002\\
$10^{13}$ & 0.164 & 0.283 & 0.063 & 0.002 & 0.140 & 0.267 & 0.035 & $<10^{-4}$\\
$10^{14}$ & 0.139 & 0.266 & 0.034 & 0 & 0.138 & 0.266 & 0.034 & 0\\
$10^{16}$ & 0.137 & 0.266 & 0.033 & 0 & 0.137 & 0.266 & 0.033 & 0\\
\bottomrule
\end{tabular}
\caption{親が合体までに超放射を受けるときの、親の合体のときのスピンの平均と、2 世代目がその値以上のスピンで生まれる割合(基準の組、$\mu=0.71\times10^{-12}\eV$)。}
\label{tab:e23boson}
\end{table}

\begin{itemize}
\item \textbf{削られるのは上の裾である}。$f\ge10^{14}\GeV$ では、親の 41\% が回転を失い、親のスピンの平均は 0.232 から 0.139、95 パーセンタイルは 0.667 から 0.262 に下がる。2 世代目の中央値は 0.675 から 0.673 でほとんど動かないが、0.72 以上は 15\% から 3\% に減り、0.78 以上はなくなる。速く回る 2 世代目は、速く回る親からしか生まれないためである。
\item \textbf{窓の $f$ でも、いちばん上の裾は減る}。$f=10^{12}\GeV$ で、回転を失った親は 17\%(古い組で 30\%)で、0.78 以上は 2.5\% から 1.5\%(0.6\%)、0.72 以上は 15\% から 13\%(11\%)になる。$f\le10^{11}\GeV$ では変わらない。
\item \textbf{親は、2 世代目より遅く回転を失う}。親の結合定数は 2 世代目(0.104)の 0.35--0.70 倍で、同じ $f$ での降下が遅い。表の変化が $f=10^{12}$--$10^{13}\GeV$ に広がっているのは、このためと、親の年齢が広く分布しているためである。
\item \textbf{仮定の効き方}: 親の合体の赤方偏移を 0.05 と 1 にすると、$f=10^{12}\GeV$ で 0.72 以上は 12\% と 14\%、0.78 以上は 1.2\% と 2.0\% である。親の質量が雲に移って減る分(211 が持ち去るとして $\alpha_g\Delta\chi$)は、平均 0.5\%、99 パーセンタイルでも 3.5\% なので、入れていない。親の $\chi_\mathrm{SR}$ の印 3--5 は、$f\le3\times10^{16}\GeV$ で 0.1\% 以下、$f=10^{17}$--$10^{18}\GeV$ で 1--3\% である。
\end{itemize}

\begin{figure}[H]
\centering
% E23: (a) 2 世代目の生まれたときのスピンの分布、(b) 生まれたときのスピンと年齢について周辺化した、0.69 以上にとどまる確率。
% データは calc/e23_natal2g.py
\begin{tikzpicture}
\begin{groupplot}[group style={group size=2 by 1, horizontal sep=17mm}, width=0.48\linewidth, height=0.44\linewidth, note]
\nextgroupplot[xmin=0.55, xmax=0.85, ymin=0, ymax=16, xlabel={生まれたときのスピン $\chi_b$}, ylabel={$dP/d\chi_b$},
  title={\small (a) 2 世代目の生まれたときのスピン},
  legend style={font=\scriptsize, at={(0.5,-0.30)}, anchor=north, fill=white, draw=none}, legend columns=1,
  legend cell align=left]
  \addplot[draw=none, fill=cM, fill opacity=0.18, forget plot] coordinates {(0.69,0) (0.69,16) (0.85,16) (0.85,0)} \closedcycle;
  \addplot[cA, line width=1.0pt] table[x=chi, y=beta_all] {figs/data/e23_hist.dat}; \addlegendentry{ボソンなし(基準の組)}
  \addplot[cA, line width=1.0pt, dashed] table[x=chi, y=beta_v100] {figs/data/e23_hist.dat}; \addlegendentry{そのうち星団に残るもの($v_\mathrm{esc}=100$ km/s)}
  \addplot[cB, line width=1.0pt] table[x=chi, y=lf12] {figs/data/e23_hist_boson.dat}; \addlegendentry{親が超放射を受ける、$f=10^{12}$ GeV}
  \addplot[cB, line width=1.0pt, dashed] table[x=chi, y=lf14] {figs/data/e23_hist_boson.dat}; \addlegendentry{親が超放射を受ける、$f=10^{14}$ GeV}
\nextgroupplot[xmode=log, xmin=1e10, xmax=1e14, ymin=0, ymax=1.05, xlabel={$f$ [GeV]}, ylabel={確率 $/$ ボソンなしの確率},
  title={\small (b) 0.69 以上にとどまる確率},
  legend style={font=\scriptsize, at={(0.5,-0.30)}, anchor=north, fill=white, draw=none}, legend columns=2,
  legend cell align=left]
  \addplot[cB, line width=1.0pt] table[x=f, y=wide] {figs/data/e23_keep.dat}; \addlegendentry{分布: 広い}
  \addplot[cM, line width=0.8pt] table[x=f, y=wide_fix078] {figs/data/e23_keep.dat}; \addlegendentry{0.78 に固定: 広い}
  \addplot[cB, line width=1.0pt, dashed] table[x=f, y=peak] {figs/data/e23_keep.dat}; \addlegendentry{分布: 星団の目安}
  \addplot[cM, line width=0.8pt, dashed] table[x=f, y=peak_fix078] {figs/data/e23_keep.dat}; \addlegendentry{0.78 に固定: 星団の目安}
  \addplot[cB, line width=1.0pt, densely dotted] table[x=f, y=slow] {figs/data/e23_keep.dat}; \addlegendentry{分布: 遅い}
  \addplot[cM, line width=0.8pt, densely dotted] table[x=f, y=slow_fix078] {figs/data/e23_keep.dat}; \addlegendentry{0.78 に固定: 遅い}
\end{groupplot}
\end{tikzpicture}

\caption{(a) 2 世代目の生まれたときのスピンの分布(基準の組)。灰色の帯は GW241011 の主星の測定の 90\% 区間(0.69 より上)。青の破線は、星団に残るものを全体に対する密度で描いた(面積は残る割合 0.13)。(b) GW241011 の主星の質量の 2 世代目が、合体のときに 0.69 以上のスピンを持つ確率を、ボソンなしの値で割ったもの($\mu=0.71\times10^{-12}\eV$)。橙は生まれたスピンと年齢の両方について周辺化した場合、灰色は生まれたスピンを 0.78 に固定した場合(E22)。}
\label{fig:e23}
\end{figure}

\subsection{生まれたスピンと年齢について周辺化した GW241011}

E22 の $P_\mathrm{keep}$ は、生まれたスピンを 0.72 か 0.78 に固定していた。ここでは式~\eqref{eq:e23chain} を最後まで通す。10 万対の親(基準の組、年齢は既定の 1 世代目)から 2 世代目を作り、2 世代目の年齢を E22 の三つの組(広い、星団の目安、遅い)から引いて、合体のときのスピンが測定の 90\% 区間の下端 0.69 以上である確率を出した。
\begin{equation}
P_\mathrm{keep}(f)=P\bigl[\chi_\mathrm{2G}\ge0.69\bigr],\qquad P_0\equiv P_\mathrm{keep}(f\to0)=P\bigl[\chi_b\ge0.69\bigr]=0.333.
\end{equation}
$P_0$ はボソンなしの値で、生まれたときに 0.69 以上である割合に等しい。ボソンの効き方は比 $P_\mathrm{keep}/P_0$ で見る(図~\ref{fig:e23}(b)、表~\ref{tab:e23f})。

\begin{table}[H]
\centering\small
\begin{tabular}{llccc}
\toprule
2 世代目の年齢 & 生まれたスピン & 0.9 & 0.5 & 0.1\\
\midrule
広い & 分布 & $7.4\times10^{10}$ & $4.2\times10^{11}$ & $3.4\times10^{12}$\\
 & 分布、残るものだけ & $5.4\times10^{10}$ & $3.4\times10^{11}$ & $3.5\times10^{12}$\\
 & 0.78 に固定 & $1.6\times10^{11}$ & $8.3\times10^{11}$ & $7.7\times10^{12}$\\
 & 0.72 に固定 & $9.7\times10^{10}$ & $4.9\times10^{11}$ & $4.6\times10^{12}$\\
\midrule
星団の目安 & 分布 & $2.6\times10^{11}$ & $9.1\times10^{11}$ & $2.8\times10^{12}$\\
 & 分布、残るものだけ & $2.0\times10^{11}$ & $7.6\times10^{11}$ & $2.7\times10^{12}$\\
 & 0.78 に固定 & $6.4\times10^{11}$ & $1.9\times10^{12}$ & $5.3\times10^{12}$\\
 & 0.72 に固定 & $3.8\times10^{11}$ & $1.1\times10^{12}$ & $3.2\times10^{12}$\\
\midrule
遅い & 分布 & $5.1\times10^{10}$ & $1.4\times10^{11}$ & $3.7\times10^{11}$\\
 & 分布、残るものだけ & $3.9\times10^{10}$ & $1.0\times10^{11}$ & $3.0\times10^{11}$\\
 & 0.78 に固定 & $1.4\times10^{11}$ & $2.4\times10^{11}$ & $6.0\times10^{11}$\\
 & 0.72 に固定 & $8.8\times10^{10}$ & $1.5\times10^{11}$ & $3.6\times10^{11}$\\
\bottomrule
\end{tabular}
\caption{$P_\mathrm{keep}/P_0$ が 0.9、0.5、0.1 になる $f$ [GeV]。「残るものだけ」は $v_\mathrm{esc}=100$ km/s($P_0=0.292$)。固定の行は $P_0=1$ で、E22 の表と同じ量である。}
\label{tab:e23f}
\end{table}

\begin{itemize}
\item \textbf{$f$ は、0.78 に固定したときの約半分になる}。比が 0.5 になる $f$ は、0.78 に固定した値の 0.49--0.57 倍、0.72 に固定した値の 0.81--0.93 倍である(三つの水準を通すと 0.36--0.62 倍、0.59--1.03 倍)。0.69 以上で生まれた 2 世代目の中央値は 0.715 で、0.69 までの余裕が小さいためである。
\item \textbf{実効的な生まれたスピンは 0.71--0.72}。同じ $f$ を与える「固定した生まれたスピン」を、0.70--0.78 に固定した計算から内挿すると、比が 0.5 の水準で 0.709--0.717、水準を変えても 0.704--0.723 である。E9 からの「現実的な側」の 0.72 は、分布で周辺化した場合に近く、わずかに楽観的だった。
\item \textbf{窓に要る年齢の比}(E22 の $R(M)$)は $f$ の上限の 2 乗に反比例するので、比が 0.5 の水準で、E22 の表の 0.72 の行の 1.15--1.54 倍、0.78 の行の 3.1--4.2 倍になる。
\item \textbf{親の超放射は、ここでは小さい}。入れない場合との差は、確率で 0.015 以下、$f$ で 0.94--1.00 倍(比が 0.5 の水準)と 0.80--0.99 倍(0.1 の水準)である。窓の $f$ では、親はまだあまり回転を失っていない。
\item \textbf{星団に残るものだけにすると、さらに厳しくなる}。$v_\mathrm{esc}=100$ km/s で $P_0=0.292$、$f$ は比が 0.5 の水準で 0.76--0.84 倍である(実効的な生まれたスピンは 0.704--0.706)。
\item \textbf{判定の閾値にはよらない}。0.69 を 0.68 と 0.70 に動かすと、$P_0$ は 0.447 と 0.250 に動くが、$f$ は 0.92--1.07 倍にしか動かない。
\item \textbf{印 3--5 の標本}は最大 5.8\% で、「保つ」と「保たない」のどちらに数えても、確率の差は 0.030 以下、$f$ は比が 0.9 と 0.5 の水準で 0.3\% 以内、0.1 の水準で最大 16\% である。
\item 生まれたスピンを固定した行は、E22 の表の値を 4\% 以内で再現する($f$ の刻みが違う)。
\end{itemize}

\paragraph{同じ $f$ で両立する度合い}
E22 は、窓を $W(M)=\max_f\min[P_\mathrm{keep}(f),P_\mathrm{drop}(f;M)]$ という一つの数にまとめた($P_\mathrm{drop}$ は、質量 $M$ の 1 世代目が、生まれたスピン 0.7 から「211 の臨界 $+0.05$」より下に落ちている確率)。$P_\mathrm{keep}$ を $P_\mathrm{keep}/P_0$ に置き換えて、生まれたスピンの分布で出し直した(\texttt{e23\_natal2g.py window}。$P_\mathrm{drop}$ は E22 の値)。

\begin{table}[H]
\centering\small
\begin{tabular}{llcccc@{\qquad}cccc}
\toprule
 & & \multicolumn{4}{c}{1 世代目: 既定} & \multicolumn{4}{c}{1 世代目: 古い}\\
2 世代目の年齢 & 生まれたスピン & $10\Msun$ & 15 & 20 & 25 & $10\Msun$ & 15 & 20 & 25\\
\midrule
広い & 分布 & 0.13 & 0.29 & 0.40 & 0.48 & 0.15 & 0.34 & 0.48 & 0.60\\
 & 分布、残るものだけ & 0.13 & 0.28 & 0.39 & 0.47 & 0.14 & 0.33 & 0.46 & 0.56\\
 & 0.72 に固定 & 0.18 & 0.33 & 0.43 & 0.50 & 0.20 & 0.38 & 0.51 & 0.63\\
 & 0.78 に固定 & 0.26 & 0.39 & 0.48 & 0.55 & 0.29 & 0.47 & 0.59 & 0.72\\
\midrule
星団の目安 & 分布 & 0.11 & 0.34 & 0.48 & 0.59 & 0.12 & 0.45 & 0.72 & 0.85\\
 & 分布、残るものだけ & 0.10 & 0.31 & 0.46 & 0.56 & 0.11 & 0.39 & 0.66 & 0.80\\
 & 0.72 に固定 & 0.13 & 0.38 & 0.53 & 0.63 & 0.16 & 0.53 & 0.78 & 0.90\\
 & 0.78 に固定 & 0.26 & 0.47 & 0.60 & 0.69 & 0.33 & 0.71 & 0.87 & 0.94\\
\midrule
遅い & 分布 & 0.00 & 0.00 & 0.07 & 0.22 & 0.00 & 0.00 & 0.08 & 0.29\\
 & 分布、残るものだけ & 0.00 & 0.00 & 0.05 & 0.16 & 0.00 & 0.00 & 0.06 & 0.19\\
 & 0.72 に固定 & 0.00 & 0.00 & 0.06 & 0.24 & 0.00 & 0.00 & 0.06 & 0.30\\
 & 0.78 に固定 & 0.00 & 0.00 & 0.19 & 0.36 & 0.00 & 0.00 & 0.24 & 0.51\\
\bottomrule
\end{tabular}
\caption{同じ $f$ で両立する度合い $W(M)$(E22 の定義で、$P_\mathrm{keep}$ の代わりに $P_\mathrm{keep}/P_0$ を使う)。固定の行は E22 の表の値にあたる($f$ の刻みの違いで、差は 0.007 以内)。}
\label{tab:e23W}
\end{table}

\begin{itemize}
\item \textbf{0.78 に固定した値より下がり、0.72 に固定した値に近い}。差は、0.78 に固定した値に対して $-0.25$ から $0.00$、0.72 に固定した値に対して $-0.07$ から $+0.02$ である。
\item \textbf{E22 のいちばん有利な組}(星団の目安の 2 世代目と古い 1 世代目)は、$20\Msun$ が 0.88 から 0.72、$15\Msun$ が 0.71 から 0.45、$10\Msun$ が 0.33 から 0.12 になる。1 世代目を既定の年齢にすると、0.48、0.34、0.11 である。軽い 1 世代目では、窓はほとんど残らない。
\item 星団に残るものだけにすると、さらに最大 0.09 下がる。
\end{itemize}

\subsection{集団解析での使い方}

\paragraph{呼び方}
式~\eqref{eq:e23chain} をそのまま書く(\texttt{share/remnant/README.md} に全文)。
\begin{verbatim}
chi_a, _ = model(m_a, chi_a0, am.sample(1, z_parent, rng), mu, f)   # 親の超放射
chi_b, _ = model(m_b, chi_b0, am.sample(1, z_parent, rng), mu, f)
out = remnant.merge(m_a, m_b, chi_a, chi_b, rng)                    # 2 世代目の誕生
chi_2g, regime = model(out["M"], out["chi"], am.sample(2, z_merge, rng), mu, f)
\end{verbatim}
\begin{itemize}
\item \textbf{親の集団は、1 世代目の集団の模型と共通にする}。親のスピンの分布(Beta の二つのパラメータ)と質量比の分布は、1 世代目の事象に当てはめるものと同じパラメータで書く。2 世代目の生まれたスピンの分布は、1 世代目の推定と一緒に動く。
\item \textbf{質量の条件づけ}: 事象の主星の質量 $M$ に対して、親の質量の和を $M/0.952$ にとる(残骸は $M$ の $\pm1\%$ に入る)。
\item \textbf{乱数}: 確率を標本の平均で出すので、$\mu$ や $f$ を変えても同じ乱数を使う(尤度がざらつかない)。
\item \textbf{手間}: \texttt{merge} は $10^5$ 対で 0.02--0.03 秒。重いのは $\chi_\mathrm{SR}$(親で 2 回、2 世代目で 1 回)なので、$(\mu,f)$ の格子の上で前もって作っておく。
\end{itemize}

\paragraph{提案}
主解析は、親のスピンを 1 世代目と共通の Beta、向きを等方、質量比を一様(下限はパラメータ)とし、親の超放射を入れ、キックでは選ばない。頑健さの節で、キックで選んだ場合($v_\mathrm{esc}=50$--300 km/s)、角度の補正、質量比の範囲を示す。GW241011 は、主星のスピンの事後分布の標本を、式~\eqref{eq:e23chain} の予言 $\chi_\mathrm{2G}$ の分布と直接比べる(0.69 の閾値は使わない)。

\begin{result}
\begin{itemize}
\item 残骸の質量・スピン・キックの部品を \texttt{share/remnant/remnant.py} に書いた。親の集団から、2 世代目の生まれたときの質量とスピンの分布を作る。
\item ボソンなしの基準の組で、2 世代目の生まれたスピンは中央値 0.675(5--95\% は 0.597--0.758)。0.69 以上は 33\%、0.72 以上は 15\%、0.78 以上は 2.5\%。
\item 親が超放射で回転を失うと、上の裾が削られる。$f\ge10^{14}\GeV$ で 0.72 以上は 3\%、0.78 以上は 0。窓の $f=10^{12}\GeV$ で 0.78 以上は 1.5\%(古い 1 世代目なら 0.6\%)。
\item 生まれたスピンと年齢について周辺化すると、GW241011 の主星が 0.69 以上にとどまる確率が半分になる $f$ は $1.4\times10^{11}$--$9.1\times10^{11}\GeV$(2 世代目の年齢の組による)。0.78 に固定した値の 0.49--0.57 倍で、0.71--0.72 で生まれた場合にあたる。
\item 同じ $f$ で両立する度合い $W$ は、星団の目安の 2 世代目と古い 1 世代目で、$20\Msun$ が 0.72、$15\Msun$ が 0.45、$10\Msun$ が 0.12(0.78 に固定した E22 の値は 0.88、0.71、0.33)。
\item キックで選ぶと $f$ はさらに 0.76--0.84 倍になるが、キックの式は粗いので、既定では選ばない。
\end{itemize}
\end{result}

\begin{caveat}
\begin{itemize}
\item \textbf{親の集団は仮定である}。0.69 以上で生まれる割合は、親が回らなければ 0、ほかの組では質量比とスピンの分布で 0.25--0.52 に動く。向きは等方(星団の中で対になった BH)に限った。向きが揃う親や、降着のある AGN 円盤の親には使えない。親は 1 世代目どうしに限り、2 世代目を親に持つ 3 世代目は入れていない。
\item \textbf{$P_0$ は式の精度と閾値に敏感である}。分布の山(0.68--0.69)が閾値 0.69 のすぐ下にあるので、スピンの当てはめの偏り(0.005)と残差(0.006)の程度の違いで、0.69 以上の割合が動く(角度の補正で 0.333 から 0.324、閾値を 0.68--0.70 に動かすと 0.447--0.250)。$f$ の値は比 $P_\mathrm{keep}/P_0$ で決めたので、この影響を受けにくい。
\item \textbf{キックの式は粗い}。球状星団の模擬計算の、親のスピン 0.5 の場合のスピンの分布(中央値 0.62)は再現できず、原因も決められなかった。キックで選んだ結果は、頑健さの確認にだけ使う。
\item \textbf{$P_\mathrm{keep}$ は説明用の量である}。閾値 0.69 は測定の 90\% 区間の下端で、集団解析では事後分布と直接比べる(D9)。
\item 親の質量の減少(平均 0.5\%)と、親の合体の赤方偏移の分布(0.3 に固定)は入れていない。
\item 数値は $\mu=0.71\times10^{-12}\eV$、$19.6\Msun$ の 2 世代目についてのもの。重い 2 世代目では親も重く、親の超放射はもっと小さい $f$ から効く。
\item $\chi_\mathrm{SR}$ は主要項の $\gamma_\mathrm{BH}$。相対論的な補正を入れると $f$ は約 1.6 倍に動く(E12)。生まれたスピンを固定した場合との比は、補正が共通なら変わらない。
\end{itemize}
\end{caveat}
